% A complete machine-checked Osterwalder-Seiler chain for the Z_2 chain.
% Lane: O (operator bridge), modules YangMills/OS/**.
% REGIME: tricotomy; NO scores, NO desk language, NO commit chronology beyond
% the reproducibility freeze.  Freeze: Lean tree of source checkpoint
% 0558dbc523337fa3d026a07fd828bb4c8a8951c2, full core build 8422 jobs,
% 2425 oracle commands (see Section 8).
\documentclass[11pt]{article}
\usepackage[margin=1.1in]{geometry}
\usepackage{amsmath,amssymb,amsthm}
\usepackage{booktabs,longtable,array}
\usepackage[colorlinks=true,linkcolor=blue,urlcolor=blue,citecolor=blue]{hyperref}

\newtheorem{theorem}{Theorem}[section]
\newtheorem{lemma}[theorem]{Lemma}
\newtheorem{proposition}[theorem]{Proposition}
\newtheorem{corollary}[theorem]{Corollary}
\newtheorem{remark}[theorem]{Remark}
\theoremstyle{definition}
\newtheorem{definition}[theorem]{Definition}

\emergencystretch=3em
\tolerance=2000

\newcommand{\lean}[1]{\texttt{#1}}
\newcommand{\anchor}{0558dbc523337fa3d026a07fd828bb4c8a8951c2}
\newcommand{\osfile}[2]{%
  \href{https://github.com/lluiseriksson/THE-ERIKSSON-PROGRAMME/blob/\anchor/YangMills/OS/#1.lean}{\nolinkurl{#2}}}
\newcommand{\osline}[3]{%
  \href{https://github.com/lluiseriksson/THE-ERIKSSON-PROGRAMME/blob/\anchor/YangMills/OS/#1.lean\#L#2}{\nolinkurl{#3}}}
\newcommand{\repofile}[2]{%
  \href{https://github.com/lluiseriksson/THE-ERIKSSON-PROGRAMME/blob/\anchor/#1}{\nolinkurl{#2}}}
\newcommand{\R}{\mathbb{R}}
\newcommand{\C}{\mathbb{C}}
\newcommand{\Z}{\mathbb{Z}}
\newcommand{\SU}{\mathrm{SU}}
\newcommand{\Om}{\Omega}
\newcommand{\inn}[2]{\langle #1,\, #2\rangle}
\newcommand{\Ex}{\mathbb{E}}
\newcommand{\ket}[1]{|#1\rangle}
\newcommand{\bra}[1]{\langle #1|}

\title{From the Gibbs Weight to the Spectral Gap:\\
A Complete Machine-Checked Osterwalder--Seiler Chain\\
for the $\Z_2$ Lattice Gauge Chain\\
\large The mathematics is textbook; what is verified\\
is that every interface composes}
\author{Lluis Eriksson}
\date{}

\begin{document}
\maketitle

\begin{abstract}
A mass gap is a statement about the spectrum of an operator, but a lattice
gauge theory is given as a measure.  The passage between the two --- the
Osterwalder--Seiler reconstruction --- proceeds through reflection positivity,
a Gelfand--Naimark--Segal space, a transfer operator, and an \emph{identification}
theorem asserting that expectations in the measure are matrix elements of that
operator.  This paper reports a complete formal verification of that passage,
in Lean 4 with \texttt{mathlib}, for the $\Z_2$ lattice gauge chain: from the
Boltzmann weight $e^{\beta s}$ to exponential decay of a correlation function
of the measure, with every arrow a machine-checked theorem and no arrow
assumed.  The endpoint is the exact identity
$\Ex_n[f(\sigma_0)f(\sigma_n)] = (\tanh\beta)^n$ for the sign observable, a
statement in which no operator occurs and whose proof goes through one, and
whose rate is exactly the verified spectral gap $\|T - |\Om\rangle\langle\Om|\|
\le \tanh\beta$.

We state the limits with the same precision as the results.  The system has one
$\Z_2$ variable per time slice, so its spatial slice is a point and its transfer
operator is a $2\times 2$ matrix; a lattice with spatial extent is not treated.
The gap is at fixed finite size and is \emph{not} volume-uniform.  The pairing
of this system is definite, so the Gelfand--Naimark--Segal quotient is the
identity and is therefore absent rather than closed.  The underlying
mathematics is textbook: Osterwalder--Seiler is from 1978 and the transfer
matrix of the Ising chain is older still.  The contribution claimed here is not
a new theorem but a verified composition: that the interfaces of the
reconstruction fit together with nothing hidden between them, exhibited on the
smallest system where all of them are simultaneously present.  Nothing in this
paper is a claim about $\SU(N)$, the continuum limit, or the Yang--Mills mass
gap.
\end{abstract}

\section{Introduction}

\subsection{The gap between a measure and a spectrum}

Constructive lattice gauge theory produces a measure.  The mass gap is a
property of an operator.  The bridge between them is the Osterwalder--Seiler
reconstruction \cite{OsterwalderSeiler}; the surrounding framework is
\cite{GlimmJaffe}.  Reflection positivity of the measure
gives a positive semi-definite pairing on functions of the fields in a half
space; quotienting by its null space and completing gives a Hilbert space;
translation in the reflected direction descends to a self-adjoint contraction
$T$ fixing the vacuum vector $\Om$; and the \emph{identification} theorem states
that the Euclidean expectation of an observable separated by $n$ time steps
from another is the matrix element $\inn{A\Om}{T^n B\Om}$.  Only after that
last step does exponential decay of correlations become a statement about the
spectrum of $T$, and only then can ``mass gap'' mean what it is supposed to
mean.

The identification is the step that is easiest to assume and hardest to supply.
It is where a formalisation can quietly become circular: if the measure is
\emph{defined} through the operator, the identification is a tautology and the
chain proves nothing.  This paper's central design constraint --- fixed in
advance, and verified mechanically rather than by inspection --- is that the two
sides be defined independently.

\subsection{What is proved here}

For the $\Z_2$ chain with $n{+}1$ time slices we define, using
$\exp$ and the bond sign and nothing else,
\[
  W_\beta(\sigma) \;=\; \prod_{t} e^{\beta\, s(\sigma_t,\sigma_{t+1})},
  \qquad
  s(i,j) = \begin{cases} +1 & i=j\\ -1 & i\neq j,\end{cases}
\]
the partition function $Z_n = \sum_\sigma W_\beta(\sigma)$ and the expectation
$\Ex_n$.  Separately, and without reference to any measure, we define the
operator
\[
  T_\beta \;=\; a(\beta)\,\mathbf{1} \;+\; b(\beta)\,S,
  \qquad
  a = \frac{e^{\beta}}{e^{\beta}+e^{-\beta}},\quad
  b = \frac{e^{-\beta}}{e^{\beta}+e^{-\beta}},
\]
where $S$ exchanges the two configurations.  The results are:

\begin{itemize}
\item $Z_n = 2\,(e^\beta+e^{-\beta})^n$, computed;
\item the identification $\Ex_n[\overline{A(\sigma_0)}\,B(\sigma_n)]
      = \inn{A\Om}{T_\beta^{\,n} B\Om}$ for every $n$ and every pair of
      observables;
\item the gap $\|T_\beta - |\Om\rangle\langle\Om|\| \le a-b = \tanh\beta$;
\item and the endpoint $\Ex_n[f(\sigma_0)f(\sigma_n)] = (\tanh\beta)^n$.
\end{itemize}

Verification data for all of the above are collected in Section~\ref{sec:repro}.

Reflection positivity of the same weight was verified separately
\cite{ErikssonRP} and the abstract criterion relating a gap to exponential
clustering was verified separately \cite{ErikssonOBridge}; both are consumed
here, so the chain is complete end to end within one formal development.

\subsection{What is deliberately not claimed}

We separate this out because thematic proximity is not mathematical
transfer.

\begin{enumerate}
\item \textbf{The system is the smallest one.}  One $\Z_2$ variable per time
slice.  The spatial slice is a point; $T_\beta$ is a $2\times 2$ matrix.  A
lattice with spatial extent has a transfer operator on a space growing with the
volume and is not treated here.
\item \textbf{The gap is not volume-uniform.}  It is a statement at fixed finite
size.  At this size the question of volume-uniformity does not arise, which is
exactly why the result is not evidence about the volume-uniform case.
\item \textbf{The Gelfand--Naimark--Segal quotient is absent, not closed.}  The
pairing of this system is already definite, so the quotient is the identity map.
That is an evasion of the difficulty, not a solution of it; a system with a
degenerate pairing still requires it.
\item \textbf{The mathematics is not new.}  The transfer matrix of the Ising
chain predates all of this, and the reconstruction is \cite{OsterwalderSeiler}.
What is new is that each interface is a machine-checked theorem and that they
compose without a gap.
\item \textbf{Nothing here concerns $\SU(N)$, the continuum limit, or the Clay
problem} \cite{Clay}.  No reduction from any statement in this paper to the
Yang--Mills mass gap is claimed, and none is known to the author.
\end{enumerate}

\section{The system and the measure}

\begin{definition}[Configurations and weight]
A configuration of the chain with $n$ bonds is a map
$\sigma : \{0,\dots,n\} \to \Z_2$.  With $s(i,j) = +1$ if $i=j$ and $-1$
otherwise, the bond weight is $k_\beta(i,j) = e^{\beta s(i,j)}$ and the path
weight is $W_\beta(\sigma) = \prod_{t=0}^{n-1} k_\beta(\sigma_t,\sigma_{t+1})$.
\end{definition}

In Lean these are \osline{Z2Identification}{67}{z2Sign},
\osline{Z2Identification}{76}{z2Bond} and
\osline{Z2Identification}{84}{z2PathWeight}.  The partition function
\osline{Z2Identification}{88}{z2Partition} and the expectation
\osline{Z2Identification}{99}{z2Expect} are built from these.

\begin{remark}[Non-circularity, verified mechanically]\label{rem:noncirc}
The definitions above and the two derived from them form a closed dependency
set:
\[
  \texttt{z2Expect} \to \{\texttt{z2PathSum},\texttt{z2Partition}\}
  \to \texttt{z2PathWeight} \to \texttt{z2Bond} \to \{\exp,\texttt{z2Sign}\}
  \to \varnothing .
\]
Printing all six definitions and inspecting the printouts shows that no
transfer operator, and neither of the coefficients $a$, $b$, occurs anywhere on
the measure side.  Because the closure is closed, this is a complete check
rather than a sample.  This is the condition that makes the identification of
Section~\ref{sec:ident} a theorem rather than a restatement of a definition.
\end{remark}

\section{Reflection positivity}

Reflection positivity of the $\Z_2$ Wilson weight was established in a
companion development \cite{ErikssonRP} and is recalled here only to place the
present chain.  The route avoids both Bochner's theorem and the Schur product
theorem by formulating the hypothesis as \emph{non-negative combination of
characters} rather than as non-negativity of Fourier coefficients; that class is
closed under products by indexing over pairs, which is what the reflected
pairing requires.  For $\Z_2$ at $\beta \ge 0$ the coefficients are
$(e^{\beta} \pm e^{-\beta})/2$, and their non-negativity is monotonicity of
$\exp$.  The relevant artifacts are
\osline{ReflectionSplitting}{131}{osPairing\_nonneg},
\osline{WilsonCharCombo}{169}{z2Wilson\_isPSDKernel} and
\osline{Z2Endpoint}{116}{z2Wilson\_reflectionPositive}, with non-degeneracy at
\osline{Z2Endpoint}{123}{z2WilsonSystem\_nondegenerate}.

\begin{remark}
For $\Z_N$ with $N > 2$ the corresponding coefficients are discrete
Bessel-type sums and their non-negativity is \emph{not} proved in this
development.  The chain of this paper is therefore specific to $\Z_2$ and does
not extend to $\Z_N$ without that missing input.
\end{remark}

\section{The transfer operator, read off the Gibbs weight}

\begin{definition}
On $\C^2$ with the Euclidean inner product let $S$ be the exchange of the two
configurations and
$T_\beta = a(\beta)\mathbf{1} + b(\beta) S$ with $a,b$ as above.
\end{definition}

The point of this definition is what it does \emph{not} contain: no eigenvector,
no projection, no spectral datum.  The coefficients are the normalised
Boltzmann weights.  In Lean, \osline{Z2Transfer}{158}{z2TransferOp}.

\begin{theorem}[The vacuum is fixed, as a consequence]\label{thm:fix}
$T_\beta \Om = \Om$ for the uniform vector $\Om = 2^{-1/2}(1,1)$.
\end{theorem}

\begin{proof}
$T_\beta \Om = (a+b)\Om$ and $a + b = 1$ by construction of the normalisation.
\end{proof}

Formally \osline{Z2Transfer}{189}{z2TransferOp\_fix}.  We stress that this is a
consequence of normalisation rather than a hypothesis; a development in which
$T\Om = \Om$ is assumed has moved the content elsewhere.

\begin{theorem}[Gap]\label{thm:gap}
For $\beta \ge 0$,
$\bigl\| T_\beta - \ket{\Om}\bra{\Om} \bigr\| \le a(\beta) - b(\beta)
= \tanh\beta < 1$,
where $\ket{\Om}\bra{\Om}$ denotes the orthogonal projection onto $\C\Om$.
\end{theorem}

\begin{proof}[Proof, as formalised]
The projected operator acts componentwise as $(a-b)$ times the projection off
the vacuum; this identity is proved from $a+b=1$ alone.  The norm bound then
follows from an explicit operator-norm estimate, with no appeal to a spectral
theorem.  The evaluation $a - b = \tanh\beta$ is proved rather than asserted.
\end{proof}

Artifacts: \osline{Z2Transfer}{261}{z2TransferOp\_gap},
\osline{Z2Transfer}{142}{z2A\_sub\_z2B\_eq\_tanh}, and the packaging
\osline{Z2Transfer}{220}{z2TransferOp\_vacuumTransfer}.

\section{The identification}\label{sec:ident}

Write $A\Om$ for the vector with entries $A(i)\,\Om(i)$.

\begin{theorem}[Identification]\label{thm:ident}
For every $n \in \mathbb{N}$ and all observables $A, B : \Z_2 \to \C$,
\[
  \Ex_n\bigl[\overline{A(\sigma_0)}\;B(\sigma_n)\bigr]
  \;=\; \inn{A\Om}{\,T_\beta^{\,n}\,B\Om}.
\]
\end{theorem}

\begin{proof}[Proof, as formalised]
Split the sum over configurations at the first spin.  Writing
$\sigma = (s,\tau)$, the first bond factors out of the path weight and the
inner sum over $s$ turns the observable $A$ into
$A'(i) = \sum_s \overline{A(s)}\,k_\beta(s,i)$, which is
$(e^\beta+e^{-\beta})$ times one application of $T_\beta$.  Each bond therefore
emits one factor of the normalisation and one application of the operator, and
closing the induction step requires only the symmetry of the kernel --- that is,
the self-adjointness of $T_\beta$.  After $n$ steps,
\[
  \sum_\sigma \overline{A(\sigma_0)}B(\sigma_n) W_\beta(\sigma)
  = (e^\beta+e^{-\beta})^n\,\inn{A}{T_\beta^{\,n} B},
\]
where $A,B$ on the right are read as vectors of $\C^2$.  Taking $A=B=1$ and
using $T_\beta \mathbf{1} = \mathbf{1}$ --- again $a+b=1$ --- gives
$Z_n = 2(e^\beta+e^{-\beta})^n$.  The two powers of the normalisation cancel in
the quotient, leaving $\tfrac12 \inn{A}{T_\beta^n B}$; and
$\inn{A\Om}{T_\beta^n B\Om} = (2^{-1/2})^2 \inn{A}{T_\beta^n B}$ is the same
number.
\end{proof}

The induction step is \osline{Z2Identification}{171}{z2PathSum\_succ}, the
partition function \osline{Z2Identification}{221}{z2Partition\_eq}, and the
theorem \osline{Z2Identification}{269}{z2\_identification}.

\begin{remark}
The cancellation of $(e^\beta+e^{-\beta})^n$ between numerator and denominator
is not arranged by hand: both sides are computed independently and the factor
appears in each.  Judge~2 of the pre-registration
(\repofile{docs/O-BRIDGE-CHARTER.md}{docs/O-BRIDGE-CHARTER.md}, Amendment~8)
required exactly this --- that the partition function be \emph{computed} rather
than symbolically cancelled --- because a development that divides by an
uncomputed $Z$ can hide an error of normalisation that would falsify the
identification.
\end{remark}

\begin{remark}[What in this argument is general, and what is not]
It is worth separating the two, because the next system will reuse one and
not the other.  The induction above --- splitting configurations at the
first spin, factoring one bond out of the path weight, and absorbing it into
the observable --- uses only that the state space is finite and that the bond
kernel is symmetric.  It does not use that the group has two elements, and it
does not use the explicit form of $k_\beta$.  Everything downstream of it
does: the dictionary between observables and vectors, the evaluation
$Z_n = 2(e^\beta+e^{-\beta})^n$, the eigenvector computation of the
next section, and the operator-norm bound of Theorem~\ref{thm:gap} are all
carried out in explicit coordinates on a two-dimensional space.  The
formalisation reflects this asymmetry in its proofs but not in its types:
every declaration is stated at the concrete state space, so the general half
of the argument is \emph{present but not abstracted}.  Porting to a larger
system therefore means re-typing the induction rather than reproving it, and
redoing the rest.
\end{remark}

\section{The endpoint}

Let $f(0) = +1$, $f(1) = -1$.

\begin{theorem}[Exact two-point function]\label{thm:endpoint}
For every $n$,
$\;\Ex_n\bigl[f(\sigma_0) f(\sigma_n)\bigr] = (\tanh\beta)^n$.
Moreover $\Ex_n[f(\sigma_0)] = 0$, so the connected two-point function
coincides with the full one.
\end{theorem}

\begin{proof}
$f$, read as a vector, satisfies $T_\beta f = (a-b) f$, so
$T_\beta^n f = (a-b)^n f$; Theorem~\ref{thm:ident} and
$\inn{f}{f} = 2$ give the value, and $a - b = \tanh\beta$.  The one-point
function is $\inn{f\Om}{\Om}$, which vanishes because $f$ sums to zero.
\end{proof}

Artifacts: \osline{Z2Identification}{346}{z2\_two\_point\_tanh},
\osline{Z2Identification}{364}{z2\_measure\_clustering} (the same in norm form),
and non-vacuity at
\osline{Z2Identification}{375}{z2\_two\_point\_ne\_zero}.

This is the statement the whole chain exists to produce.  Its left-hand side
contains no operator: it is an expectation in a measure defined from
$e^{\beta s}$.  Its right-hand side is a number that came out of the spectrum of
$T_\beta$.  Its rate is exactly the gap of Theorem~\ref{thm:gap}.  That the
three agree is the content of the reconstruction, and here it is checked rather
than asserted.

\begin{remark}[On what an exact answer does and does not show]
For this system the two-point function can be computed directly, without any of
the machinery above; the value $(\tanh\beta)^n$ is classical \cite{Simon}.  The verification
reported here is therefore not evidence that the machinery is \emph{necessary}
--- at this size it plainly is not.  It is evidence that the machinery is
\emph{correct}: an independently constructed operator-theoretic route reproduces
the classical answer exactly, which is the check one wants before applying the
same route where no direct computation is available.
\end{remark}

\section{The abstract criterion, consumed}

The companion development \cite{ErikssonOBridge} proves, for a self-adjoint
contraction $T$ with $T\Om = \Om$ and $\|\Om\| = 1$, that a bound
$\|T - \ket{\Om}\bra{\Om}\| \le r$ is \emph{equivalent} to exponential decay of
the connected two-point function at rate $r$ along a family with dense span,
with an unconstrained constant per observable.  Theorem~\ref{thm:gap} supplies a
hypothesis of that criterion from an operator that came from the measure, so the
criterion is consumed rather than merely stated:
\osline{SharpBridge}{218}{connCorrC\_le\_of\_gap} and
\osline{SharpBridge}{270}{sharp\_clustering\_iff\_gap}, resting on
\osline{DenseClustering}{225}{norm\_le\_of\_dense\_decayDomain}.

\begin{remark}
Before the present work, every witness of that criterion in the development had
the form $r\mathbf{1} + (1-r)P$ --- an operator built \emph{from} the conclusion.
An adversarial audit and an external review both identified this as the
criterion's outstanding weakness.  $T_\beta$ is not of that form; it is read off
the Gibbs weight.  We record that the objection applied, and that it is answered
for this system and for no larger one.
\end{remark}

\section{Formalisation architecture}

The development is Lean~4 \cite{Lean4} with \texttt{mathlib} \cite{Mathlib}
pinned to an exact commit.  Three design decisions are worth recording because
they are what made the chain composable.

\paragraph{Separation of the two sides into disjoint definition sets.}
The measure side and the operator side share no definitions; they meet in
exactly one lemma, which states that the normalised bond weight equals the
transfer matrix entry.  This is what makes Remark~\ref{rem:noncirc} checkable by
printing definitions.  A development that shares a ``kernel'' definition between
the two sides cannot make that check.

\paragraph{Observables as plain functions, vectors as a separate type.}
Observables are $\Z_2 \to \C$; vectors live in the Euclidean space.  A single
explicit dictionary maps between them.  Identifying the two types would have
made the identification theorem read as a triviality.

\paragraph{Induction on the path, not on the operator.}
The proof of Theorem~\ref{thm:ident} inducts on the number of bonds, splitting
configurations with the standard cons equivalence, and never manipulates
$T^n$ symbolically.  The operator power is assembled by the induction rather
than assumed to interact with the sum.

\paragraph{House notes.}  Two recurrences cost builds and are recorded for
re-use.  First, goals produced by the self-adjointness/symmetry interface sit
over the \emph{linear-map} coercion, so a lemma stated about the continuous
linear map does not apply until the coercion is bridged.  Second, the two
unfolding lemmas for powers act on opposite sides, and pairing an iteration
lemma that unfolds on the left with a power lemma that unfolds on the right
leaves the two sides of an induction step misaligned in a way whose error
message points nowhere near the cause.

\section{Reproducibility and verification data}\label{sec:repro}

All artifacts are at source checkpoint \texttt{\anchor} on branch \texttt{main}
of \url{https://github.com/lluiseriksson/THE-ERIKSSON-PROGRAMME}.  Every link in
this paper was resolved against that commit before compilation.

\begin{center}
\begin{tabular}{ll}
\toprule
Quantity & Value \\
\midrule
Full core build & 8422 jobs, success \\
Oracle commands & 2425 \\
\quad with axiom dependencies & 2403 \\
\quad axiom-free & 22 \\
Distinct axioms across all outputs & \texttt{propext}, \texttt{Quot.sound}, \texttt{Classical.choice} \\
Occurrences of \texttt{sorryAx} & 0 \\
Project axioms & 0 \\
Errors & 0 \\
\bottomrule
\end{tabular}
\end{center}

The oracle script is \repofile{oracle\_check.lean}{oracle\_check.lean}; the core
root is \repofile{YangMillsCore.lean}{YangMillsCore.lean}.  The five criteria this work was
required to meet --- including the non-circularity condition of
Remark~\ref{rem:noncirc} and the computed partition function of
Section~\ref{sec:ident} --- were fixed in
\repofile{docs/O-BRIDGE-CHARTER.md}{docs/O-BRIDGE-CHARTER.md}, Amendment~8,
in a commit preceding the existence of the module they govern; their verdicts
are recorded in
\repofile{docs/VERIFICATION-LEDGER.md}{docs/VERIFICATION-LEDGER.md},
Addendum~512.

\section{Theorem--artifact map}

\begin{center}
\begin{longtable}{p{0.30\textwidth} p{0.62\textwidth}}
\toprule
Statement & Artifact \\
\midrule
\endhead
Bond weight, path weight & \osline{Z2Identification}{76}{z2Bond},
  \osline{Z2Identification}{84}{z2PathWeight} \\
Expectation in the measure & \osline{Z2Identification}{99}{z2Expect} \\
Transfer operator & \osline{Z2Transfer}{158}{z2TransferOp} \\
Theorem~\ref{thm:fix} (vacuum fixed) & \osline{Z2Transfer}{189}{z2TransferOp\_fix} \\
Transfer data & \osline{Z2Transfer}{220}{z2TransferOp\_vacuumTransfer} \\
Theorem~\ref{thm:gap} (gap) & \osline{Z2Transfer}{261}{z2TransferOp\_gap} \\
Rate equals $\tanh\beta$ & \osline{Z2Transfer}{142}{z2A\_sub\_z2B\_eq\_tanh} \\
Induction step & \osline{Z2Identification}{171}{z2PathSum\_succ} \\
Partition function & \osline{Z2Identification}{221}{z2Partition\_eq} \\
Theorem~\ref{thm:ident} (identification) & \osline{Z2Identification}{269}{z2\_identification} \\
Theorem~\ref{thm:endpoint} (endpoint) & \osline{Z2Identification}{346}{z2\_two\_point\_tanh} \\
Endpoint, norm form & \osline{Z2Identification}{364}{z2\_measure\_clustering} \\
Non-vacuity & \osline{Z2Identification}{375}{z2\_two\_point\_ne\_zero} \\
Reflection positivity & \osline{Z2Endpoint}{116}{z2Wilson\_reflectionPositive} \\
Non-degeneracy & \osline{Z2Endpoint}{123}{z2WilsonSystem\_nondegenerate} \\
Positive-definite kernel & \osline{WilsonCharCombo}{169}{z2Wilson\_isPSDKernel} \\
Reflected pairing & \osline{ReflectionSplitting}{131}{osPairing\_nonneg} \\
Gap $\Rightarrow$ clustering & \osline{SharpBridge}{218}{connCorrC\_le\_of\_gap} \\
Clustering $\Leftrightarrow$ gap & \osline{SharpBridge}{270}{sharp\_clustering\_iff\_gap} \\
Dense-family criterion & \osline{DenseClustering}{225}{norm\_le\_of\_dense\_decayDomain} \\
\bottomrule
\end{longtable}
\end{center}

\section{Continuation}

The next step is not a larger gauge group but a harder pairing.  The
Gelfand--Naimark--Segal quotient does no work in this paper because the pairing
is definite; the first genuine test is a system whose pairing is degenerate, so
that the quotient must be constructed and its properties transported through it.
After that comes spatial extent --- a transfer operator acting on a space that
grows with the volume --- and only after that does volume-uniformity become a
meaningful question.  Volume-uniformity is the first point at which any of this
could bear on a mass gap, and it is not reached here.  We state this ordering
because the temptation in this area is to jump to the last item; the intervening
steps are where constructions of this kind are expected to fail, and the present
one has not yet been asked the question.

\begin{thebibliography}{9}

\bibitem{OsterwalderSeiler}
K.~Osterwalder and E.~Seiler,
\emph{Gauge field theories on a lattice},
Annals of Physics \textbf{110} (1978), 440--471.

\bibitem{GlimmJaffe}
J.~Glimm and A.~Jaffe,
\emph{Quantum Physics: A Functional Integral Point of View},
2nd ed., Springer, 1987.

\bibitem{Simon}
B.~Simon,
\emph{The Statistical Mechanics of Lattice Gases, Volume I},
Princeton University Press, 1993.

\bibitem{ErikssonOBridge}
L.~Eriksson,
\emph{Clustering and the Transfer-Operator Gap: A Machine-Checked Dense-Family
Criterion}, viXra, 2026.

\bibitem{ErikssonRP}
L.~Eriksson,
\emph{Reflection Positivity for Abelian Lattice Gauge Theory, Machine-Checked},
viXra, 2026.

\bibitem{Lean4}
L.~de~Moura and S.~Ullrich,
\emph{The Lean 4 theorem prover and programming language},
CADE-28, Lecture Notes in Computer Science \textbf{12699} (2021), 625--635.

\bibitem{Mathlib}
The mathlib Community,
\emph{The Lean mathematical library},
CPP 2020, 367--381.

\bibitem{Clay}
A.~Jaffe and E.~Witten,
\emph{Quantum Yang--Mills theory},
Clay Mathematics Institute Millennium Prize Problem description, 2000.

\end{thebibliography}

\end{document}
